
\documentclass{article}
\usepackage{colortbl}
\usepackage{makecell}
\usepackage{multirow}
\usepackage{supertabular}

\begin{document}

\newcounter{utterance}

\centering \large Interaction Transcript for game `hot\_air\_balloon', experiment `air\_balloon\_survival\_en\_reasoning off\_easy', episode 1 with Qwen3{-}32B{-}t0.0.
\vspace{24pt}

{ \footnotesize  \setcounter{utterance}{1}
\setlength{\tabcolsep}{0pt}
\begin{supertabular}{c@{$\;$}|p{.15\linewidth}@{}p{.15\linewidth}p{.15\linewidth}p{.15\linewidth}p{.15\linewidth}p{.15\linewidth}}
    \# & $\;$A & \multicolumn{4}{c}{Game Master} & $\;\:$B\\
    \hline

    \theutterance \stepcounter{utterance}  
    & & \multicolumn{4}{p{0.6\linewidth}}{
        \cellcolor[rgb]{0.9,0.9,0.9}{
            \makecell[{{p{\linewidth}}}]{
                \texttt{\tiny{[P1$\langle$GM]}}
                \texttt{You are participating in a collaborative negotiation game.} \\
\\ 
\texttt{Together with another participant, you must agree on a single set of items that will be kept. Each of you has your own view of how much each item matters to you (importance). You do not know how the other participant values the items. Additionally, you are given the effort each item demands.} \\
\texttt{You may only agree on a set if the total effort of the selected items does not exceed a shared limit:} \\
\\ 
\texttt{LIMIT = 4100} \\
\\ 
\texttt{Here are the individual item effort values:} \\
\\ 
\texttt{Item effort = \{"A19": 108, "B91": 133, "A30": 669, "B11": 835, "C47": 460, "C75": 537, "C26": 573, "C73": 737, "A47": 865, "A15": 855, "A75": 596, "A39": 719, "A56": 533, "C24": 549, "B26": 31\}} \\
\\ 
\texttt{Here is your personal view on the importance of each item:} \\
\\ 
\texttt{Item importance values = \{"A19": 138, "B91": 583, "A30": 822, "B11": 783, "C47": 65, "C75": 262, "C26": 121, "C73": 508, "A47": 780, "A15": 461, "A75": 484, "A39": 668, "A56": 389, "C24": 808, "B26": 215\}} \\
\\ 
\texttt{Goal:} \\
\\ 
\texttt{Your goal is to negotiate a shared set of items that benefits you as much as possible (i.e., maximizes total importance to YOU), while staying within the LIMIT. You are not required to make a PROPOSAL in every message {-} you can simply negotiate as well. All tactics are allowed!} \\
\\ 
\texttt{Interaction Protocol:} \\
\\ 
\texttt{You may only use the following structured formats in a message:} \\
\\ 
\texttt{PROPOSAL: \{'A', 'B', 'C', …\}} \\
\texttt{Propose keeping exactly those items.} \\
\texttt{REFUSE: \{'A', 'B', 'C', …\}} \\
\texttt{Explicitly reject opponent's proposal.} \\
\texttt{ARGUMENT: \{'...'\}} \\
\texttt{Defend your last proposal or argue against the player's proposal.} \\
\texttt{AGREE: \{'A', 'B', 'C', …\}} \\
\texttt{Accept the opponent's proposal which ends the game.} \\
\\ 
\\ 
\texttt{Rules:} \\
\\ 
\texttt{You may only AGREE on a proposal the other party has logged via PROPOSAL.} \\
\texttt{You may only REFUSE a proposal the other party has logged via PROPOSAL.} \\
\texttt{Total effort of any PROPOSAL or AGREE set must be ≤ LIMIT.} \\
\texttt{Do NOT reveal your hidden importance scores.} \\
\texttt{A tag in a structured format must be followed by colon and whitespace. The argument must be a python set containing 0 or more strings.} \\
\texttt{So, it must be of the form TAG: \{...\}} \\
\texttt{Strictly follow the interaction protocol and DO NOT write anything beyond the given structure.} \\
\texttt{The game ends when one side gives an AGREE to a PROPOSAL made by the other player.} \\
\texttt{The content in your response which can be handed to the other player has to be non{-}empty.} \\
\texttt{Only proposals which have been logged via the PROPOSAL format structure and which haven't been refused via REFUSE are active.} \\
\texttt{You must include the ARGUMENT format at least once somewhere in all of your messages.} \\
            }
        }
    }
    & \\ \\

    \theutterance \stepcounter{utterance}  
    & \multicolumn{4}{p{0.6\linewidth}}{
        \cellcolor[rgb]{0.9,0.9,0.9}{
            \makecell[{{p{\linewidth}}}]{
                \texttt{\tiny{[P1$\rangle$GM]}}
                \texttt{<think>} \\
\texttt{Okay, let's see. I need to negotiate a set of items that maximizes my total importance without exceeding the effort limit of 4100. First, I should figure out which items are most important to me and check their combined effort.} \\
\\ 
\texttt{Looking at my importance values: A30 is 822, B91 is 583, A47 is 780, A15 is 461, A39 is 668, A56 is 389, A75 is 484, A19 is 138, C24 is 808, C73 is 508, B11 is 783, C47 is 65, C26 is 121, C75 is 262, B26 is 215. } \\
\\ 
\texttt{The highest importance items are A30 (822), B91 (583), A47 (780), B11 (783), C24 (808). Let me check their efforts. A30's effort is 669, B91 is 133, A47 is 865, B11 is 835, C24 is 549. Wait, adding those up: 669 + 133 = 802, plus 865 is 1667, plus 835 is 2502, plus 549 is 3051. That's under 4100. But maybe I can add more items. Let's see. Next high importance is A39 (668) with effort 719. Adding that would make total 3051 + 719 = 3770. Then maybe A15 (461) at 855 effort. 3770 + 855 = 4625, which exceeds the limit. So can't add that. How about C24 is already in. What about C73 (508) with effort 737? 3051 + 737 = 3788. Still} \\
            }
        }
    }
    & & \\ \\

    \theutterance \stepcounter{utterance}  
    & & \multicolumn{4}{p{0.6\linewidth}}{
        \cellcolor[rgb]{0.9,0.9,0.9}{
            \makecell[{{p{\linewidth}}}]{
                \texttt{\tiny{[P1$\langle$GM]}}
                \texttt{Your response did not contain an argument. You must include the structured format ARGUMENT: \{...\} somewhere in your response. Try again} \\
            }
        }
    }
    & \\ \\

    \theutterance \stepcounter{utterance}  
    & \multicolumn{4}{p{0.6\linewidth}}{
        \cellcolor[rgb]{0.9,0.9,0.9}{
            \makecell[{{p{\linewidth}}}]{
                \texttt{\tiny{[P1$\rangle$GM]}}
                \texttt{<think>} \\
\texttt{Okay, let me try again. I need to make sure I include an argument in my response. The user wants me to propose a set of items that maximizes my importance without exceeding the effort limit. Let me go through the items again.} \\
\\ 
\texttt{My top importance items are A30 (822), B91 (583), A47 (780), B11 (783), C24 (808). Let's check their efforts: A30 is 669, B91 is 133, A47 is 865, B11 is 835, C24 is 549. Adding those up: 669 + 133 = 802, plus 865 is 1667, plus 835 is 2502, plus 549 is 3051. That's under 4100. Now, can I add more? Next is A39 (668) with effort 719. 3051 + 719 = 3770. Still under. Then maybe C73 (508) at 737 effort. 3770 + 737 = 4507, which is over 4100. So can't add that. How about C26 (121) with effort 573? 3051 + 573 = 3624. Still under. But C26's importance is low. Maybe better to add A56 (389) at 533 effort. 3051 + 533 = 3584. Then A75 (484) at 596: 3584 + 596 = 4180, which is over. So maybe stop there. So the set would be A30, B91, A47, B11, C24, A39, A56. Let me check the total effort again: 669 + 133 = 802, +865=1667, +835=25} \\
            }
        }
    }
    & & \\ \\

    \theutterance \stepcounter{utterance}  
    & & \multicolumn{4}{p{0.6\linewidth}}{
        \cellcolor[rgb]{0.9,0.9,0.9}{
            \makecell[{{p{\linewidth}}}]{
                \texttt{\tiny{[P1$\langle$GM]}}
                \texttt{Your response did not contain an argument. You must include the structured format ARGUMENT: \{...\} somewhere in your response. Try again} \\
            }
        }
    }
    & \\ \\

    \theutterance \stepcounter{utterance}  
    & \multicolumn{4}{p{0.6\linewidth}}{
        \cellcolor[rgb]{0.9,0.9,0.9}{
            \makecell[{{p{\linewidth}}}]{
                \texttt{\tiny{[P1$\rangle$GM]}}
                \texttt{<think>} \\
\texttt{Okay, I need to make sure I include an argument in my response. Let me start by calculating the total effort for the items I proposed. The items are A30, B91, A47, B11, C24, A39, and A56. Let me add up their efforts:} \\
\\ 
\texttt{{-} A30: 669} \\
\texttt{{-} B91: 133} \\
\texttt{{-} A47: 865} \\
\texttt{{-} B11: 835} \\
\texttt{{-} C24: 549} \\
\texttt{{-} A39: 719} \\
\texttt{{-} A56: 533} \\
\\ 
\texttt{Adding these together: 669 + 133 = 802; 802 + 865 = 1667; 1667 + 835 = 2502; 2502 + 549 = 3051; 3051 + 719 = 3770; 3770 + 533 = 4303. Wait, that's over 4100. Oh no, I made a mistake here. The total effort is 4303, which exceeds the limit. That's a problem. I need to adjust the set to stay under 4100.} \\
\\ 
\texttt{Let me remove the item with the lowest importance in the set to reduce the effort. The lowest importance in the current set is A56 (389). If I remove A56 (533 effort), the total becomes 4303 {-} 533 = 3770. Now, can I add another item? Let's see. The remaining items not in the set include A15 (855 effort, 461 importance), C73 (737 effort, 508 importance), and others. Adding A15 would make the total 3770 + 855 = 4625, which is way over. C73 would be 3770 + 737 = 4507, still over. How about C26 (573 effort, 121 importance)? 3770 + 573 =} \\
            }
        }
    }
    & & \\ \\

    \theutterance \stepcounter{utterance}  
    & & & \multicolumn{2}{p{0.3\linewidth}}{
        \cellcolor[rgb]{0.9,0.9,0.9}{
            \makecell[{{p{\linewidth}}}]{
                \texttt{\tiny{[GM$|$GM]}}
                \texttt{abort} \\
            }
        }
    }
    & & \\ \\

    \theutterance \stepcounter{utterance}  
    & & & \multicolumn{2}{p{0.3\linewidth}}{
        \cellcolor[rgb]{0.9,0.9,0.9}{
            \makecell[{{p{\linewidth}}}]{
                \texttt{\tiny{[GM$|$GM]}}
                \texttt{end game} \\
            }
        }
    }
    & & \\ \\

\end{supertabular}
}

\end{document}
